\section{Introduction}

Large firms tend to pay higher wages than small firms. This empirical regularity, usually known as the firm-size wage premium, is one of the classic facts in labor economics, and it often remains after conditioning on observed worker characteristics \citep{BrownMedoff1989,OiIdson1999,troske1999evidence}. Nonetheless, a firm-size wage premium conditional on observables should be interpreted carefully. It means that workers in larger firms earn more than observably similar workers in smaller firms, but it does not prove that moving a given worker to a larger firm would raise aggregate output, since unobserved worker quality or other omitted factors may still matter. At the same time, the persistence of the premium after controlling for observables is informative: it is consistent with allocative frictions even if it is not direct proof of it. Evidence from developing economies and matched employer--employee studies supports this cautious interpretation: firm-size effects often remain after stronger controls for worker heterogeneity, but their interpretation depends on the underlying source of the employer-associated wage gap \citep{SoderbomTealWambugu2005,reed2019large,AbowdKramarzMargolis1999,CardHeiningKline2013,Song2019,AlvarezBenguriaEngbomMoser2018,GerardLagosSeverniniCard2021,EngbomMoser2022}. The firm-size wage premium is therefore a diagnostic statistic rather than a causal estimate.

Our motivation to measure the firm-size wage premium in Colombia is both empirical and policy-oriented. A recent policy diagnosis argues that low-income households are disproportionately linked to own-account work, very small businesses, informality, and low-earning jobs, whereas higher-income households are more often connected to larger and more formal employers \citep{eslava2026crecimiento}. This diagnosis is suggestive, but the observed association between firm size and household income may reflect worker sorting or formality rather than firm scale itself. Nonetheless, the wage gradient associated with firm size has not been systematically documented in Colombia. Existing work has studied wage differentials by sector and firm characteristics in the formal labor market \citep{iregui2012diferenciales}, but less is known about the relationship between firm size and wages once solo work,
micro-firms, and informal employment are included. The central empirical challenge is therefore to document the firm-size wage gradient while distinguishing it from worker characteristics and the formal-informal wage gap.

This paper documents the firm-size wage premium in Colombia using harmonized microdata from the Gran Encuesta Integrada de Hogares (GEIH) for 2008--2019 and 2021--2025. The central questions are descriptive and econometric: do workers in larger firms earn higher real hourly labor income than comparable workers in solo work or smaller firms? How much of the raw gradient remains after accounting for observable worker characteristics, formality, and sector-year conditions? Has the premium widened or narrowed over time? Does the relationship between firm size and wages differ between formal and informal workers? The GEIH does not identify employers or observe marginal products, so the resulting estimates are conditional wage gaps, not causal effects of firm size. They are nevertheless informative because persistent conditional gaps in a labor market with concentrated micro-firm employment are a useful diagnostic of segmentation and possible allocative frictions.

We document five main facts. First, employment is concentrated at the bottom of the firm-size distribution: solo workers and firms with up to ten workers account for about two thirds of employment. Second, the raw wage gradient by firm size is large and widespread across sectors, education groups, gender, and formality status. Third, the gradient attenuates substantially after adding worker controls, sector-year fixed effects, and especially formality, but it does not disappear: in the full specification, workers in firms with 101 or more employees earn about 43 percent more than otherwise comparable solo workers. Fourth, the conditional premium is positive throughout the period, with no statistically conclusive evidence of either a secular decline or a systematic widening. Fifth, the formality status affects the firm-size wage premium. Formal workers earn more than informal workers within every firm-size category, but the conditional firm-size wage premium is much steeper among informal workers. Within formal employment, the broad size gradient is weak relative to formal solo workers, although workers in 101+ firms earn more than workers in intermediate-size formal firms. If the wage gaps are interpreted as a markers of allocative frictions rather than as pure sorting, the evidence points to three possible margins for efficiency gains: movement from informal to formal employment, scaling of informal micro-units, and, more selectively, movement or growth within the formal sector toward larger employers.

The rest of the paper is organized as follows. Section~2 describes the GEIH data, the harmonization of firm-size categories, and the construction of real hourly labor income. Section~3 presents the descriptive evidence on employment shares and wage gaps by firm size, formality, education, gender, sector, and year. Section~4 introduces the econometric specifications. Section~5 reports the main regression results, the dynamic estimates, and the formality interactions. Section~6 discusses interpretation and limitations, Section~7 draws policy implications, and Section~8 concludes.
